\documentclass[12font]{article}
\title{\textbf{Correlation Functions in a Finite-Size Kitaev Model}}
\author{Chen-Chih Wang  }
\date{\textit{Department of Physics, National Tsing Hua University, Hsinchu 30013, Taiwan} \\ (Dated: \today)}
\usepackage[margin=2cm]{geometry}
%\usepackage[fleqn]{amsmath}
%\usepackage{CJK,amsmath}
\usepackage{amsfonts,amssymb}
\usepackage{fancyhdr}
%\pagestyle{fancy}
\usepackage{textcomp}
\usepackage{graphicx}
\usepackage{float}
\usepackage{color}
\usepackage{tikz}
\usepackage{multicol}
\usepackage{threeparttable}
\usepackage{amsthm}
\usepackage{mathtools}
\usepackage{hyperref}
\urlstyle{same}
\usepackage{caption}
\usepackage{subcaption}
\usepackage{dsfont}
\usepackage{qcircuit} % quantum circuit

\newcommand{\vac}{\mrm{vac}}
\newcommand{\Bra}[1]{\left\langle #1 \right|}
\newcommand{\Ket}[1]{\left| #1 \right\rangle}
\newcommand{\bra}[1]{\langle #1 |}
\newcommand{\ket}[1]{| #1 \rangle}
\newcommand{\anglelr}[1]{\left\langle #1 \right\rangle}
\newcommand{\braket}[2]{\langle #1 | #2 \rangle}
\newcommand{\Bbraket}[2]{\left\langle #1 \right|\left. #2 \right\rangle}
\newcommand{\dif}[1]{\frac{\mathrm{d}}{\mathrm{d}#1}}
\newcommand{\diff}[2]{\frac{\mathrm{d}#1}{\mathrm{d}#2}}
\newcommand{\gd}{\boldsymbol{\nabla}}
\newcommand{\dive}{\boldsymbol{\nabla}\cdot}
\newcommand{\curl}{\boldsymbol{\nabla}\times}
\newcommand{\lrno}[1]{\left.#1\right.}
\newcommand{\lr}[1]{\left(  #1  \right)}
\newcommand{\lrm}[1]{\left[  #1  \right]}
\newcommand{\lrg}[1]{\left\{  #1  \right\}}
\newcommand{\abs}[1]{\left|  #1  \right|}
\newcommand{\de}{\partial}
\newcommand{\pdif}[1]{\frac{\partial}{\partial#1}}
\newcommand{\pdiff}[2]{\frac{\partial#1}{\partial#2}}
\newcommand{\mrm}[1]{\mathrm{#1}}
\newcommand{\bo}[1]{\mathbf{#1}}
\newcommand{\bog}[1]{\boldsymbol{#1}}
\newcommand{\vint}[3]{\int_{#1}#2\cdot\mathrm{d}\mathbf{#3}}
\newcommand{\voint}[3]{\oint_{#1}#2\cdot\mathrm{d}\mathbf{#3}}
\newcommand{\Vint}{\int\mrm{d}^3\bo{r}'}
\newcommand{\rr}{\mathbf{r} - \mathbf{r}'}
\newcommand{\arr}{\left|\mathbf{r} - \mathbf{r}'\right|}
\newcommand{\nx}{\hat{\bo{n}}_1}
\newcommand{\ny}{\hat{\bo{n}}_2}
\newcommand{\alert}[1]{{\color{red}#1}}
\newcommand{\alertV}[1]{{\color{violet}#1}}
\newcommand{\Tr}{\mrm{Tr}}
\newcommand{\tr}{\mrm{tr}}
\newcommand{\norm}[1]{\left\lVert#1\right\rVert}
\newcommand{\Qcirc}[1]{\begin{minipage}[h]{0.7cm}
	\vspace{0pt}
	\Qcircuit @C=1em @R=.7em {#1}
   \end{minipage}}
\newcommand{\dimer}[2]{\sigma_{#1}^{\alpha_{#1#2}}\sigma_{#2}^{\alpha_{#1#2}}}

\begin{document}
\maketitle

%\fontsize{12pt}{18pt}\selectfont

\begin{abstract}
    In this note, we discuss how to compute multi-spin correlation functions in a finite-size Kitaev model, which only contains four distinct plaquettes. These correlation functions are crucial in constructing the density matrix of the Kitaev ground state.
\end{abstract}

\section{Foundations}

\subsection{Correlation Functions}
As mentioned in our work, at the ground state, the multi-spin correlation function of the Kitaev model is finite if and only if it can be written in the form of a string. Therefore, we only need to discuss the following form of the correlation function
\[
\anglelr{
    \prod_{\anglelr{ij}\in\mathfrak{D}}
    \dimer{i}{j}
}
\]
This kind of correlation function can be exactly evaluated in the Majorana representation. Substituting the transformation 
\[
\dimer{i}{j} \to i b^{\alpha_{ij}}_i c_i  i b^{\alpha_{ij}}_j c_j = \hat{u}_{ij} ic_jc_i
\]
into the correlation function and taking the trivial gauge, we have
\[
\anglelr{
    \prod_{\anglelr{ij}\in\mathfrak{D}}
    \dimer{i}{j}
}
=
\anglelr{
    \prod_{\anglelr{ij}\in\mathfrak{D}}
    ic_jc_i
}
\]

\subsection{Minimal Size for Non-trivial Topology}
For numerical efficiency, we tend to choose the smallest possible system. However, the system must still be sufficiently large such that the topological order still exists. A requirement for this is that the ground state manifold should possess four independent topological sectors generated by global Wilson loops if the geometry of the system is a torus.

We start from the Wen plaquette model. There are two requirements for a Wen plaquette system to possess four topologically degenerate ground states:
\begin{itemize}
    \item The lattice must contain an even number of rows or columns in both the horizontal and vertical directions.
    
    \item A global Wilson loop cannot be generated or destroyed by local plaquette operators.
\end{itemize}
The physical meaning of the first requirement is to preserve the checkerboard pattern of the lattice. If there are adjacent plaquettes (cells) that are of the same type (black or white, in a checkerboard sense), then the two types of plaquettes become identical, and the original two different Wilson loops become the same. This breakdown of the checkerboard pattern reduces the original fourfold topological ground state degeneracy down to twofold, which is not adequate to describe the physics of a toric code system.

According to the two requirements, \textbf{the smallest possible topologically non-trivial Wen plaquette model contains eight ($8 = 2\times 4$) distinct plaquettes and eight sites}. A $2\times 2$ checkerboard is not adequate for the system to have four topological sectors since it violates the second requirement\footnote{The \textit{global Wilson loop} is identical to a plaquette operator}. 

\section{Finite-Size Correlation}
We only have to focus on the two-point correlation functions. Suppose that the system has a periodic boundary condition, for which the 


\end{document}